\section{Introduction}\label{sec:introduction}

Multipoint Seshadri constants measure how much positivity an ample line
bundle retains after imposing simultaneous conditions at several points.
Their volume bound depends only on the top self-intersection of the line
bundle and the number of points. The question is whether, for sufficiently
many general points, any curve can impose a stricter bound. We prove that
no such curve exists at a geometric generic tuple in positive
characteristic.

Let $k$ be an algebraically closed field of characteristic $p>0$, let
$X/k$ be a smooth integral projective variety of dimension $n$, and let
$L$ be an ample line bundle. For $r\ge1$, write
\[
 U_r=\{(x_1,\ldots,x_r)\in X^r:x_i\ne x_j\text{ for }i\ne j\}
\]
for the open configuration space of ordered distinct points. Put
$K_r=\overline{k(U_r)}$, and let $p_1,\ldots,p_r\in X(K_r)$ be the
points specified by its geometric generic point. Define
\begin{equation}\label{eq:generic-seshadri}
 \varepsilon_r(X,L)=\inf_C
 \frac{L_{K_r}\cdot C}{\sum_{i=1}^r\mult_{p_i}C},
\end{equation}
where $C$ ranges over integral curves in $X_{K_r}$ meeting at least one
$p_i$. Multiplicity at a point outside $C$ is zero. Throughout, we use
this ordinary, curve-defined Seshadri constant.

\begin{theorem}\label{thm:main}
Let $X/k$ be a smooth integral projective variety over an algebraically
closed field of characteristic $p>0$, with $\dim X=n\ge3$, and let $L$
be ample. There is an integer $r_0=r_0(X,L)$ such that
\[
 \varepsilon_r(X,L)=\left(\frac{L^n}{r}\right)^{1/n}
 \qquad\text{for every integer }r\ge r_0.
\]
\end{theorem}

The same conclusion holds when $n=2$; see Remark~\ref{rem:surfaces}.

The geometric generic formulation includes countable algebraically closed
fields. It also gives a precise meaning to general position without
choosing points in an infinite intersection of open subsets. The equality
is exact for each sufficiently large point count. On the blow-up of the
geometric generic points, it says that
\[
 \pi^*L_{K_r}-(L^n/r)^{1/n}\sum_{i=1}^r E_i
\]
is a nef real divisor class with top self-intersection zero.

\subsection{History and the characteristic issue}

Nagata's work on Hilbert's fourteenth problem led to his conjecture on
the degrees of plane curves with assigned multiplicities at general
points, and to its proof when the number of points is a square at least $16$
\cite{Nagata1959}. The qualitative Nagata--Biran--Szemberg assertion
extends eventual maximality to polarized varieties. Ro\'e and Ross
formulate it in arbitrary dimension over uncountable algebraically closed
fields and distinguish the ordinary
Seshadri constant from its higher-dimensional cycle analogues
\cite[Conjecture~1.3 and Remark~1.2]{RoeRoss2009}.
Theorem~\ref{thm:main} proves its positive-characteristic form, with
general position expressed by the geometric generic tuple, in the stated
dimensions.

Several preceding results explain the volume scale of this question.
Biran's packing stability theorem gives a symplectic analogue on
closed four-manifolds with rational symplectic class \cite{Biran1999}.
K\"uchle obtains asymptotically optimal lower bounds for multipoint
Seshadri constants on complex projective varieties \cite{Kuechle1996}.
Product inequalities compare
point counts on a variety with those on projective space
\cite[Theorem~1.1 and Remark~1.2]{RoeRoss2009}.
These are distinct from exact algebraic maximality for every sufficiently
large integer $r$.

Our proof adapts the graded-support construction of the companion
\emph{Maximal Multipoint Seshadri Constants in Higher Dimensions}
\cite[Sections~2--6]{OpenAI2026Higher}, which treats complex varieties.
Its two-coordinate reshaping follows the nodal exponent bounds and
primitive lattice directions in the surface companion
\cite[Propositions~3.2 and~3.4 and Lemma~4.3]{OpenAI2026Surfaces}.
We prove every required statement here. The passage to positive
characteristic requires an algebraic replacement for analytic limits and
a nodal equation that retains two distinct branches in characteristic
two. Formal substitutions over $k[[s]]$ supply the first; the curve
$y^2+xy-x^3=0$ supplies the second.

The relation between positivity and asymptotic jet separation is
classical \cite[Theorem~6.4]{Demailly1992}. The use of leading monomials
belongs to the valuation approach to linear series
\cite{Okounkov1996,LazarsfeldMustata2009,KavehKhovanskii2012};
Ito gives finite-degree comparisons between multipoint jet separation
and initial monomial spaces \cite[Proposition~5.7 and Theorem~5.8]{Ito2013}.
Park and Shin give characteristic-independent Okounkov-body descriptions
of one-point Seshadri constants for complete series of nef and big
divisors \cite[Theorem~1.2(1)]{ParkShin2021}.
Monomial interpolation also underlies Dumnicki's diagram-cutting method
\cite[Proposition~13 and Theorem~14]{Dumnicki2007}.
The statements below isolate two useful forms of these ideas: formal
transfer of arbitrary finite jets, and reshaping of a graded family with
prescribed real simplex intercepts. The latter keeps both the exact
cardinality in every degree and multiplication between degrees. A short
counting argument then turns full asymptotic density into exact filling
of compact subsets of the interior.

\subsection{Proof strategy}

Put $V=L^n$ and, for a fixed point count $r$, put $w=(V/r)^{1/n}$.
For a line bundle $A$ at a smooth point $x$, its jets through order $m$
form $A_x/\mathfrak m_x^{m+1}A_x$. Separating those jets at a tuple
means that sections surject onto the direct sum of these spaces. Such
separation by $H^0(X,L^{\otimes N})$ gives the curve inequality
\[
 NL\cdot C\ge m\sum_i\mult_{p_i}C.
\]
It therefore suffices to obtain jets through order $\lfloor N\theta\rfloor$
for every fixed $0<\theta<w$ and all sufficiently large $N$.

We encode sections by finite sets $S_N\subset\Z^n$ of leading exponents,
with the multiplication property
\[
 S_N+S_M\subseteq S_{N+M}.
\]
Their monomial spaces on $\T^n=(\mathbb G_{m,k})^n$ are
$M(S_N)=\Span_k\{z^\alpha:\alpha\in S_N\}$.
Section~\ref{sec:transfer} shows how jet separation for such monomials
can be transferred to the original functions or sections. The transfer
is a finite-degree rank argument. Its auxiliary weights may depend on
$N$, while the orders defining the graded sets remain fixed.

For positive intercepts $a_1,\ldots,a_n$, write
\[
 \Delta(a_1,\ldots,a_n)
 =\left\{x\in\R_{\ge0}^n:\sum_i x_i/a_i\le1\right\}.
\]
A complete-intersection flag in Section~\ref{sec:flag} gives
\[
 S_N\subseteq N\Delta(1/D,\ldots,1/D,D^{n-1}V),
 \qquad \#S_N=h^0(X,L^{\otimes N}),
\]
where $DL$ is very ample. The simplex has volume $V/n!$, equal to the
leading coefficient of the section count.

The main geometric step, Section~\ref{sec:reshape}, replaces two
intercepts $A,B$ by $w,AB/w$ for any sufficiently small prescribed real
$w>0$. It preserves cardinalities, the graded property, and the
implication from jet separation after the replacement to jet separation
before it. For a suitable leading-term order, a node bounds the exponents
by a quadrilateral of the same area. Integral changes of coordinates and successive compressions,
which move interval slices to start at zero, turn this quadrilateral
into the desired triangle. The other coordinates serve only as labels.

For large $r$, repeating this step yields graded supports $F_N$ inside
\[
 NP_*,\qquad P_* =\Delta(w,\ldots,w,rw),
 \qquad \#F_N=\frac{V}{n!}N^n+o(N^n).
\]
Section~\ref{sec:interpolation} proves that every lattice point in a
fixed compact interior region belongs to $F_N$ in large degree. Indeed,
each such point has on the order of $N^n$ decompositions into two
nearly equal degrees, whereas the missing points exclude only $o(N^n)$
of them. The long final intercept of $P_*$ then permits polynomial
interpolation at $r$ points differing only in their last coordinate.
This interpolation uses the Chinese remainder theorem and is valid in
every characteristic.

Each transferred jet-separation statement holds on a nonempty open
subset of $U_r$, hence at its geometric generic point.
Section~\ref{sec:positivity} converts these statements into the required
lower bound for every curve and proves the matching volume upper bound
on the blow-up.
